% Source: https://artofproblemsolving.com/community/c3h2476712_rectangular_prism_find_the_space_diagonal
% Stable identity: aops-topic-2476712-q01
Let the edge lengths be integers $a\le b\le c$. Then
\[abc=720,\qquad 2(ab+bc+ca)=484\iff ab+bc+ca=242.\]
Dividing the second equation by the first,
\[\dfrac1a+\dfrac1b+\dfrac1c=\dfrac{242}{720}=\dfrac{121}{360}.\]
Since $a$ is the smallest edge, $\dfrac3a\ge\dfrac{121}{360}$, so $a\le\dfrac{1080}{121}<9$; and $\dfrac1a<\dfrac{121}{360}$, so $a>\dfrac{360}{121}>2$. Thus $3\le a\le8$, and $a$ divides $720$, so $a\in\{3,4,5,6,8\}$. For each $a$, $bc=\dfrac{720}a$ and $b+c=\dfrac{242-bc}{a}$:
\begin{itemize}
\item $a=3$: $bc=240$, $b+c=\dfrac23$; \quad $a=4$: $bc=180$, $b+c=\dfrac{62}4$;
\item $a=5$: $bc=144$, $b+c=\dfrac{98}5$; \quad $a=6$: $bc=120$, $b+c=\dfrac{122}6$;
\item $a=8$: $bc=90$, $b+c=19$.
\end{itemize}
Only $a=8$ gives an integer $b+c$, and then $b,c$ are the roots of $t^2-19t+90=(t-9)(t-10)$. So the edges are $8,9,10$ (check: $8\cdot9\cdot10=720$ and $72+90+80=242$).

The longest segment joining two vertices is a space diagonal. By the Pythagorean theorem, a face diagonal is $\sqrt{a^2+b^2}$, and it is perpendicular to the third edge, so the space diagonal is
\[\sqrt{a^2+b^2+c^2}=\sqrt{(a+b+c)^2-2(ab+bc+ca)}=\sqrt{27^2-484}=\sqrt{245}=\boxed{7\sqrt5\text{ cm}}.\]

\begin{asy}
usepackage("amsmath");
size(6cm);
defaultpen(fontsize(10pt));
real a=10,c=8; pair dv=(5,2.8);
pair P0=(0,0),P1=(a,0),P2=(a,c),P3=(0,c);
pair Q0=P0+dv,Q1=P1+dv,Q2=P2+dv,Q3=P3+dv;
draw(P0--P1--P2--P3--cycle,linewidth(0.8));
draw(P1--Q1--Q2--P2,linewidth(0.8)); draw(P3--Q3--Q2,linewidth(0.8));
draw(P0--Q0,dashed); draw(Q0--Q1,dashed); draw(Q0--Q3,dashed);
draw(P0--Q1,blue+dashed);
draw(P0--Q2,red+linewidth(1));
dot(P0); dot(Q1); dot(Q2);
label("$10$",(P0+P1)/2,S); label("$8$",(Q1+Q2)/2,E); label("$9$",(P1+Q1)/2,SE);
label("$\sqrt{10^2+9^2}$",P0+0.45*(Q1-P0),dir(100),blue+fontsize(8pt));
label("$7\sqrt5$",P0+0.45*(Q2-P0),dir(-60),red+fontsize(9pt));
\end{asy}
